\documentclass[12pt]{article}
\usepackage[utf8]{inputenc}
\usepackage[portuguese]{babel}
\usepackage{amsmath, amssymb}
\usepackage{graphicx}
\usepackage[margin=2.5cm]{geometry}

\title{1ª Avaliação — Processamento Digital de Sinais}
\author{Milena Freire Britto Neves, Ana Iara Loayza Costa, José Victor Brito Costa}
\date{28/07/2026}

\begin{document}
\maketitle

\section*{Questão 1}

\subsection*{a) }

Um sinal senoidal discreto $x[n] = A\cos(\omega_0 n + \theta)$ é periódico se, e somente se,
existir $N \in \mathbb{Z}^+$ e $r \in \mathbb{Z}$ tais que:
\[
\omega_0 N = 2\pi r
\]

Para o sinal dado, $\omega_0 = \dfrac{12\pi}{5}$ e $\theta = 0$. Substituindo:
\[
\frac{12\pi}{5}N = 2\pi r \;\Rightarrow\; N = \frac{5r}{6}
\]

O menor valor inteiro de $r$ que torna $N$ inteiro é $r = 6$, resultando em:
\[
N = \frac{5 \cdot 6}{6} = 5
\]

Portanto, $x[n]$ é periódico, com período fundamental $N = 5$.

\bigskip
\noindent\textit{Gráfico comprovando a periodicidade:}

\begin{figure}[h!]
    \centering
    \includegraphics[width=0.7\textwidth]{images/grafico_1a.png}
\end{figure}

\subsection*{b) }

\textbf{Multiplicando por 2:} $\omega_0' = \dfrac{24\pi}{5}$
\[
\frac{24\pi}{5}N = 2\pi r \;\Rightarrow\; N = \frac{5r}{12}
\]
Menor $r$ inteiro válido: $r = 12 \Rightarrow N = 5$.

 O período permanece $N = 5$, mesmo dobrando a frequência. Isso explica uma
propriedade importante dos sinais discretos: diferentemente do caso contínuo
($T = 2\pi/f$, inversamente proporcional à frequência), no caso discreto essa
relação não é garantida.

\textbf{Multiplicando por 5:} $\omega_0'' = 12\pi$
\[
12\pi N = 2\pi r \;\Rightarrow\; N = \frac{r}{6}
\]
Menor $r$ inteiro válido: $r = 6 \Rightarrow N = 1$. Como $12\pi n$ é sempre múltiplo inteiro de $2\pi$, temos $\cos(12\pi n) = 1$
para todo $n$, ou seja, o sinal se torna constante (período trivial $N=1$).

\begin{figure}[h!]
    \centering
    \includegraphics[width=0.7\textwidth]{images/grafico_1b_x2.png}
\end{figure}

\begin{figure}[h!]
    \centering
    \includegraphics[width=0.7\textwidth]{images/grafico_1b_x5.png}
\end{figure}

\section*{Questão 3}

Para um sistema linear e invariante no tempo, a saída é dada pela convolução
entre o sinal de entrada $x[n]$ e a resposta ao impulso $h[n]$:
\[
y[n]=(x*h)[n]=\sum_{k=-\infty}^{\infty}x[k]h[n-k].
\]

Como o enunciado estabelece que ambos os sinais são causais,
$x[n]=h[n]=0$ para $n<0$. Assim, o somatório pode ser limitado a
\[
y[n]=\sum_{k=0}^{n}x[k]h[n-k].
\]

Se $x$ possui $N_x$ amostras e $h$ possui $N_h$ amostras, a convolução completa
possui
\[
N_y=N_x+N_h-1
\]
amostras. Na implementação, cada valor de $y[n]$ é inicialmente igual a zero.
Para cada par de índices válidos $k$ e $m$, o produto $x[k]h[m]$ contribui para
a amostra de índice $n=k+m$. Logo:
\[
y[k+m]\leftarrow y[k+m]+x[k]h[m].
\]
Esse procedimento implementa diretamente a definição da convolução, sem usar
funções prontas.

\noindent A função em Python é:
\begin{verbatim}
import numpy as np

def Convolucao(x, h):
    x = np.asarray(x)
    h = np.asarray(h)

    if x.ndim != 1 or h.ndim != 1:
        raise ValueError("x e h devem ser vetores unidimensionais.")
    if x.size == 0 or h.size == 0:
        raise ValueError("x e h não podem ser vetores vazios.")

    tipo = np.result_type(x.dtype, h.dtype)
    y = np.zeros(x.size + h.size - 1, dtype=tipo)

    for k in range(x.size):
        for m in range(h.size):
            y[k + m] += x[k] * h[m]

    return y
\end{verbatim}

Como exemplo, sejam
\[
x[n]=\{1,2,3\}
\qquad\text{e}\qquad
h[n]=\left\{1,\frac{1}{2}\right\},
\]
com a primeira amostra de cada vetor correspondendo a $n=0$. Aplicando a
definição:
\[
\begin{aligned}
y[0]&=x[0]h[0]=1,\\
y[1]&=x[0]h[1]+x[1]h[0]=\frac{1}{2}+2=\frac{5}{2},\\
y[2]&=x[1]h[1]+x[2]h[0]=1+3=4,\\
y[3]&=x[2]h[1]=\frac{3}{2}.
\end{aligned}
\]
Portanto,
\[
\boxed{y[n]=\left\{1,\frac{5}{2},4,\frac{3}{2}\right\}}.
\]

\section*{Questão 4}

A equação fornecida pode ser reescrita como
\[
y[n]=(1+r)y[n-1]+x[n].
\]
Definindo $a=1+r$, obtém-se o sistema de primeira ordem
\[
y[n]-ay[n-1]=x[n].
\]

Como o aporte mensal é constante e o primeiro aporte ocorre ao final do primeiro
mês, sua sequência é $x[n]=x\,u[n-1]$. A resposta ao impulso do sistema é
\[
H(z)=\frac{1}{1-az^{-1}}
\qquad\Longleftrightarrow\qquad
h[n]=a^n u[n].
\]

A resposta total é a soma da resposta à condição inicial com a resposta de estado
nulo. Para um saldo inicial $y[0]=y_0$, a primeira parcela é
\[
y_{\mathrm{inicial}}[n]=a^n y_0.
\]
Já a parcela causada pelos aportes pode ser calculada por convolução:
\[
\begin{aligned}
y_{\mathrm{aportes}}[n]
  &=\sum_{k=1}^{n}x\,h[n-k] \\
  &=x\sum_{k=1}^{n}a^{n-k}
   =x\sum_{m=0}^{n-1}a^m.
\end{aligned}
\]
Para $r\neq 0$, essa soma geométrica resulta em
\[
y_{\mathrm{aportes}}[n]
=x\frac{a^n-1}{a-1}
=x\frac{(1+r)^n-1}{r}.
\]
Portanto, a solução fechada, sem iteração recursiva, é
\[
\boxed{
y[n]=y_0(1+r)^n+x\frac{(1+r)^n-1}{r}
},\qquad r\neq 0.
\]
Quando $r=0$, o limite da expressão fornece
\[
\boxed{y[n]=y_0+nx}.
\]

\noindent A função em Python correspondente é:
\begin{verbatim}
import numpy as np

def Investimento(y0, r, x, n):
    if not isinstance(n, (int, np.integer)) or isinstance(n, (bool, np.bool_)):
        raise TypeError("n deve ser um número inteiro de meses.")
    if n < 0:
        raise ValueError("n não pode ser negativo.")
    if r <= -1:
        raise ValueError("r deve ser maior que -1.")

    if np.isclose(r, 0.0):
        return y0 + n*x

    fator = (1 + r)**n
    return y0*fator + x*np.expm1(n*np.log1p(r))/r
\end{verbatim}

Por exemplo, para $y_0=1000$, rentabilidade mensal de $1\%$
($r=0{,}01$), aporte mensal de $200$ e prazo de $12$ meses:
\[
\begin{aligned}
\mathrm{Investimento}(1000,0{,}01,200,12)
&=1000(1{,}01)^{12}
  +200\frac{(1{,}01)^{12}-1}{0{,}01}\\
&\approx 3663{,}33.
\end{aligned}
\]

\end{document}
